% ---------------------------------------------------------------------
% 1.5  ESCOLHA DO FORMATO DE PONTO FIXO          -> nova seção da Metodologia
% ---------------------------------------------------------------------
% ORIGEM: TCC/Análise bits.md (resumo da autora), com as correções de 2026-11-04:
%   nomes 1A-3D no lugar das letras A-P; IoTh; "condição de teste" no lugar de
%   "cenário"; critério de escolha lexicográfico (não a razão perda/ocupação);
%   5 formatos finais (não 3); AP_TRN_ZERO corrigido; tabela de velocidades
%   retirada (já está em FUNDAMENTACAO/modelo-de-canal.tex).

\subsection{\texorpdfstring{SNR increase}{SNR increase}}

To select the $\langle W, I\rangle$ format (Section~\ref{subsubsec:ponto-fixo}), the BCD-MMSE algorithm was simulated using both fixed-point and 64-bit floating-point matrix inversion. Only the inversion of $\mathbf G_c$ differs between the two simulations; however, the SNR and SER are measured for the entire uplink. Each format was evaluated based on the SNR increase—that is, the additional SNR required by the uplink to achieve an $\mathrm{SER} = 10^{-2}$ compared to the floating-point case. This increase corresponds to the horizontal distance between the two SER curves, which is the standard method for comparing two schemes \cite[p.~55]{tse-viswanath}.

Two thresholds were employed. A 0.1~dB threshold represents a near-ideal format, where the SER curve is virtually identical to that of the floating-point implementation. A 3~dB threshold serves as the maximum acceptable increase. An increase of 3~dB is equivalent to doubling the power \cite[p.~73]{tse-viswanath}. Consequently, at this threshold, each UE would need to transmit with twice the energy per symbol to compensate for the error introduced by the matrix inversion (Section~\ref{subsec:snr}).



\subsection{\texorpdfstring{\textcolor{red}{Varredura inicial}}{Varredura inicial}}
\textcolor{red}{A primeira etapa variou a parte inteira e a parte fracionária separadamente, com a outra parte fixa (Tabela~\ref{tab:varredura-inicial}). Cada tamanho foi testado nos 35 modos de quantização e de estouro (Tabela~\ref{tab:modos-apfixed}), com o BCD-MMSE em 1 e em 4 iterações.}


\begin{table}[htbp]
\caption{\textcolor{red}{Tamanhos testados na varredura inicial.}}
\label{tab:varredura-inicial}
\centering
\footnotesize
\begin{tabular}{lll}
\toprule
Varredura & Bits inteiros ($I$) & Bits fracionários ($F$) \\
\midrule
Parte inteira & 6, 8, 10, 12, 14, 16, 20, 24, 32 & 32 \\
Parte fracionária & 14 & 4, 6, 8, 10, 12, 16, 20, 24, 32 \\
\bottomrule
\end{tabular}
\legend{\textcolor{red}{Fonte: a autora.}}
\end{table}
% NOTA: 9 + 9 = 18 tamanhos, mas <46,14> (I = 14, F = 32) aparece nas duas
%   varreduras: 17 tamanhos por varredura de configuração e condição.

\subsection{\texorpdfstring{\textcolor{red}{Configurações de teste para as matrizes}}{Configurações de teste para as matrizes}}
\textcolor{red}{Na varredura inicial, foram usadas as 16 configurações cuja curva em ponto flutuante atinge $\mathrm{SER} = 10^{-2}$ entre $-10$ e 10~dB: 1A--1H, 2A, 2B, 2E, 2F, 3A, 3B, 3E e 3F. Rodar as 16 nas 5 condições de teste daria 80 casos, cada um com 17 tamanhos e 35 modos. Por isso, foram escolhidas só 5 configurações.}
% CONFERIDO (2026-11-06): metrica_ponto_operacao.csv tem os 80 casos, pontos de operação de
%   -10,0 a 2,36 dB. 1C (configuração 5) no (192; 3,5 GHz) cruza no limite: SER = 1,00e-2 em
%   -10 dB. As QPSK com M = 256 (2C, 2D, 2G, 2H, 3C, 3D, 3G, 3H) cruzam abaixo de -10 dB
%   (curvas_qpsk_256, grade de -20 a 10 dB); a autora preferiu não citar isso por enquanto.
% ANTIGO (2026-11-06, texto reduzido; tirada a frase "As configurações QPSK dos grupos 2 e 3 foram incluídas depois"):
% \textcolor{red}{Na varredura inicial, havia 16 configurações: 1A--1H, 2A, 2B, 2E, 2F, 3A, 3B, 3E e 3F. As configurações QPSK dos grupos 2 e 3 foram incluídas depois. Simular as 16 configurações nas 5 condições de teste daria 80 varreduras, cada uma com 17 tamanhos e todos os modos. Por isso, a varredura inicial usou só 5 configurações.}
% ANTIGO (2026-11-06): o parágrafo começava em "Simular as 16 configurações ..." sem dizer quais eram.
% FONTE: software/testes/mapa_cenarios_v4.py, LETRAS (numeração antiga 1-16 = 1A-1H, 2A, 2E,
%   2B, 2F, 3A, 3E, 3B, 3F); a autora incluiu as QPSK dos grupos 2 e 3 por último.

\textcolor{red}{Os 80 casos foram ordenados pelo maior valor absoluto entre as partes reais e imaginárias de $\mathbf{G}_c$, pois o circuito guarda as duas partes em separado. Esse valor é chamado aqui de maior $|$parte$|$ de $\mathbf{G}_c$. Foram usadas as matrizes dos dois pontos de SNR em volta do ponto de operação, $\mathrm{SER} = 10^{-2}$ em ponto flutuante. Em cada caso, foi tomado o percentil 99, para que poucas matrizes raras não decidam a ordem.}
% ANTIGO (2026-11-06, texto reduzido; juntado com o parágrafo do percentil):
% \textcolor{red}{Para escolher essas configurações, os 80 casos (16 configurações nas 5 condições de teste) foram ordenados pelo tamanho dos elementos de $\mathbf{G}_c$. Em cada caso, o ponto de operação é a SNR em que a curva em ponto flutuante atinge $\mathrm{SER} = 10^{-2}$. Foram usadas as matrizes $\mathbf{G}_c$ dos dois pontos de SNR simulados em volta desse ponto. Os elementos de $\mathbf{G}_c$ são complexos, e o circuito guarda a parte real e a parte imaginária em separado. Por isso, de cada matriz foi tomado o maior valor absoluto entre as partes reais e imaginárias de todos os seus elementos. Esse valor é chamado aqui de maior $|$parte$|$ de $\mathbf{G}_c$.}

% (parágrafo do percentil juntado ao anterior)
% ANTIGO (2026-11-06, texto reduzido):
% \textcolor{red}{Cada caso tem 1600 ou 3200 matrizes. Para resumir esses valores em um número, foi usado o percentil 99: o valor que fica acima de 99\% dos valores do caso. Assim, só 1\% das matrizes tem um elemento maior que ele. O percentil foi usado no lugar do máximo para que poucas matrizes raras não decidam a ordem.}
% FONTE: software/testes/metrica_ponto_operacao.py (docstring, passos 1-4) e
%   resultados/4_varredura_formatos/metrica_ponto_operacao.csv (80 linhas, coluna gc_p99).
%   Ponto de operação: referência em ponto flutuante [1], BCD-MMSE com 1 iteração.
%   Matrizes por caso: 100 realizações x C DUs x 2 SNRs = 1600 (C = 8) ou 3200 (C = 16).
%   As configurações QPSK do grupo 1 (5 a 8) usam as matrizes das 16-QAM (1 a 4) nas suas
%   próprias SNRs: a modulação não muda G_c.
% "Partes em separado": hls_inverse / cholesky_ops.h usam ap_fixed para real e imaginária
%   (CONFERIR a forma de citar; a definição do formato está em FUNDAMENTACAO/inversao.tex).
% JUSTIFICATIVA do percentil 99 (no lugar do máximo): escrita por mim, não está registrada
%   pela autora. CONFERIR se é esse o motivo.

\textcolor{red}{As configurações 1G, 1H, 2E e 2F ocupam as 17 primeiras posições. Elas têm os maiores elementos e, por isso, precisam de mais bits inteiros. A configuração 1A tem os menores elementos e foi usada como caso base (Tabela~\ref{tab:gc-ponto-operacao}).}
% ANTIGO (2026-11-06, texto reduzido):
% \textcolor{red}{Quatro delas são as configurações em que a matriz $\mathbf{G}_c$, no ponto de operação $\mathrm{SER} = 10^{-2}$, tem os elementos de maior valor absoluto: 2E, 2F, 1G e 1H. Elementos maiores precisam de mais bits inteiros para não saturar a entrada da inversão. Assim, essas são as configurações mais exigentes. A quinta é a configuração 1A, que tem os elementos de menor valor absoluto. Ela foi usada como caso base, para comparação. A Tabela~\ref{tab:gc-ponto-operacao} mostra esses valores.}

% (parágrafo das posições juntado ao anterior)
% ANTIGO (2026-11-06, texto reduzido):
% \textcolor{red}{As quatro configurações exigentes aparecem repetidas no topo da ordenação. As configurações 1G e 1H ocupam as 10 primeiras posições, com todas as suas condições de teste. As configurações 2E e 2F ocupam as posições 11 a 17. Já a configuração 1A ocupa a última posição, e as suas outras condições ficam entre as posições 75 e 78.}
% FONTE: metrica_ponto_operacao.csv, coluna posicao. 1G = 6, 1H = 8: posições 1-10.
%   2E = 10, 2F = 12: posições 11-17 (as condições restantes em 20, 27 e 29; N = 24 nas
%   duas últimas). 1A = 1: posições 75-78 e 80 (79 é 1B com N = 24).
%   Logo depois de 2E e 2F vêm 3F (18) e 3E (19), com valores parecidos (948 e 944).

\begin{table}[htbp]
\caption{\textcolor{red}{SNR no ponto de operação e maior elemento de $\mathbf{G}_c$ nas configurações da varredura inicial.}}
\label{tab:gc-ponto-operacao}
\centering
\footnotesize
\begin{tabular}{ccccc}
\toprule
 & \multicolumn{2}{c}{SNR no ponto de operação (dB)} & \multicolumn{2}{c}{Maior $|$parte$|$ de $\mathbf{G}_c$ (percentil 99)} \\
\cmidrule(lr){2-3}\cmidrule(lr){4-5}
Configuração & $N = 192$; 3,5 GHz & $N = 168$; 28 GHz & $N = 192$; 3,5 GHz & $N = 168$; 28 GHz \\
\midrule
1A & $-2{,}6$ & $-2{,}4$ & 63   & 66   \\
\textcolor{red}{1G} & $-9{,}3$ & $-9{,}1$ & 2266 & 2383 \\
\textcolor{red}{1H} & $-8{,}7$ & $-8{,}6$ & 2470 & 2385 \\
\textcolor{red}{2E} & $-5{,}8$ & $-5{,}6$ & 962  & 999  \\
\textcolor{red}{2F} & $-5{,}6$ & $-5{,}5$ & 983  & 987  \\
\bottomrule
\end{tabular}
\legend{\textcolor{red}{Fonte: a autora.}}
% ANTIGO (legenda): "Parte" é a parte real ou imaginária de cada elemento. -> foi para o texto.
\end{table}
% CORREÇÃO (2026-11-04): as linhas estavam com os nomes trocados (2E/2F com os valores de
%   1G/1H e vice-versa). Pelo metrica_ponto_operacao.csv: config 6 (1G) = -9,29/-9,13 dB e
%   2265,6/2382,9; config 8 (1H) = -8,74/-8,59 e 2470,0/2385,2; config 10 (2E) = -5,80/-5,60
%   e 961,5/999,0; config 12 (2F) = -5,61/-5,49 e 983,4/986,5. As QPSK (1G, 1H) operam em
%   SNR mais baixa e por isso têm G_c maior. Ordem antiga das linhas: 1A, 2E, 2F, 1G, 1H.
% CORREÇÃO antiga de nomes (mapa LETRAS de software/testes/mapa_cenarios_v4.py):
%   A = config 1 -> 1A; M = 10 -> 2E; N = 12 -> 2F; P = 6 -> 1G; O = 8 -> 1H.
%   No Análise bits.md, a primeira tabela troca o K de O e P (O com K = 8, P com
%   K = 12). Pelos números, 6 (1G) tem K = 8 e 8 (1H) tem K = 12; a tabela do
%   grupo 16x16, mais abaixo no mesmo arquivo, está certa. As linhas de 1G e 1H
%   acima seguem os números (P = 1G, O = 1H). CONFERIR os valores no CSV de origem.
% A CONFERIR: a tabela mostra só as 5 escolhidas. Para sustentar "as 4 maiores",
%   vale mostrar (ou citar em nota) o percentil 99 das 16 configurações.
% Ressalva (memória do projeto): o máximo absoluto de |G_c| chegou a 3960 em 1G/1H
%   a -10 dB, contra o limite de 4096 de I = 13; por isso I = 12 satura.

% SEPARAÇÃO (2026-11-04): a subseção dos modos de quantização e de estouro (texto e
%   Tabela tab:modos-apfixed) foi para FUNDAMENTACAO/inversao.tex.

\subsection{\texorpdfstring{\textcolor{red}{Síntese e seleção dos formatos para a validação}}{Síntese e seleção dos formatos para a validação}}
\textcolor{red}{Os tamanhos e modos da varredura inicial também foram sintetizados no Vitis HLS, para medir a ocupação de recursos (LUT, DSP, FF e BRAM). A ocupação de cada circuito é a do recurso mais usado no FPGA. Para cada parte da palavra, foram escolhidos os tamanhos com menor ocupação em duas classes: aumento de SNR de até 3~dB e de até 0,1~dB. Combinando os tamanhos escolhidos para a parte inteira e para a fracionária, ficaram 10 formatos para a etapa de validação.}

\textcolor{red}{A escolha foi feita para cada par (configuração, condição de teste), separada por limite de aumento de SNR e por número de iterações. Quando dois tamanhos empatam na ocupação, fica o que se repete em mais pares. O formato de um par junta a parte inteira e a parte fracionária escolhidas para ele. Os 10 formatos são todos os que apareceram nessa escolha. A Tabela~\ref{tab:formatos-origem} mostra em quais configurações cada formato foi escolhido.}

\begin{table}[htbp]
\caption{\textcolor{red}{Configurações em que cada formato foi escolhido na varredura inicial.}}
\label{tab:formatos-origem}
\centering
\footnotesize
\begin{tabular}{ccccc}
\toprule
 & \multicolumn{2}{c}{1 iteração} & \multicolumn{2}{c}{4 iterações} \\
\cmidrule(lr){2-3}\cmidrule(lr){4-5}
Formato & Aumento até 3~dB & Aumento até 0,1~dB & Aumento até 3~dB & Aumento até 0,1~dB \\
\midrule
$\langle 14,6\rangle$  & 1A         & --                       & --             & --                       \\
$\langle 16,6\rangle$  & --         & --                       & 1A             & --                       \\
$\langle 18,8\rangle$  & --         & 1A                       & --             & --                       \\
$\langle 20,8\rangle$  & --         & --                       & --             & 1A                       \\
$\langle 22,12\rangle$ & 1H, 2E, 2F & --                       & --             & --                       \\
$\langle 24,12\rangle$ & 1G         & 2E, 2F                   & 1G, 1H, 2E, 2F & --                       \\
$\langle 26,12\rangle$ & --         & 1G$^{a}$, 1H$^{a}$       & --             & 2E, 2F                   \\
$\langle 28,12\rangle$ & --         & --                       & --             & 1G$^{a}$, 1H$^{a}$       \\
$\langle 28,14\rangle$ & --         & 1G$^{b}$, 1H$^{b}$       & --             & --                       \\
$\langle 30,14\rangle$ & --         & --                       & --             & 1G$^{b}$, 1H$^{b}$       \\
\bottomrule
\end{tabular}
\legend{\textcolor{red}{Fonte: a autora. Sem marca: nas 5 condições de teste. $^{a}$Só na TDLB100-400. $^{b}$Nas outras 4 condições.}}
\end{table}
% FONTE: resultados/4_varredura_formatos/v2_casos_extremos/restricoes_varredura_4_detalhe.csv
%   (restricoes_v2.py). Cada linha do CSV é um lado (I ou F) de um par; o formato do par
%   é <I + F, I>. Nomes: 1 = 1A, 6 = 1G, 8 = 1H, 10 = 2E, 12 = 2F. TDLA30-10 = TDLA30-5.
% DESEMPATE (restricoes_v2.py): 1) a opção (tamanho e modo) que empata em mais células;
%   2) no lado de F, AP_RND_ZERO; 3) menor perda; 4) ordem dos modos. O texto cita só o 1.
% A "repetição" só desempata. Nenhum formato que apareceu foi descartado: os 10 são a
%   união das escolhas.

\textcolor{red}{Os formatos menores vêm do caso base, a configuração 1A. Os formatos com 12 e 14 bits inteiros vêm das configurações mais exigentes. Os formatos com mais bits fracionários vêm do limite de 0,1~dB e de 4 iterações.}

% CORREÇÃO: o Análise bits.md fala em minimizar "aumento de SNR / % de ocupação".
%   O critério aplicado é em ordem: primeiro o limite de perda, depois a menor
%   ocupação. Não é uma razão.
% A CONFERIR: como os pares (I, F) foram combinados nos 10 formatos.

\subsection{\texorpdfstring{\textcolor{red}{Validação}}{Validação}}
\textcolor{red}{A varredura inicial fixou uma parte da palavra enquanto variava a outra, e usou só 5 configurações. As duas coisas podem mudar o recurso mais usado e o melhor modo. Por isso, os 10 formatos foram simulados de novo, com os 35 modos, nas \textcolor{red}{24} configurações e nas 5 condições de teste, com o BCD-MMSE em 1 iteração.}
% ANTIGO: "nas 16 configurações". As 8 configurações QPSK com M = 256 (2C, 2D, 2G, 2H, 3C,
%   3D, 3G, 3H) entraram na [4] v5, completa (1750 rodadas por configuração = 10 x 35 x 5).
% FONTE: software/testes/LEIAME.md, varredura [4] v4 (28.000 rodadas = 10 formatos x
%   35 modos x 16 configurações x 5 condições), só nas SNRs em volta de SER = 1e-2.

\textcolor{red}{O circuito tem um modo só. Por isso, cada formato recebeu um modo fixo \textcolor{red}{em cada grupo de configurações (Seção~\ref{sec:configuracoes}). Todas as configurações de um grupo usam o mesmo circuito.} O modo foi escolhido pelos seguintes critérios, nesta ordem:}
% ANTIGO: "cada formato recebeu um modo fixo para os circuitos 16 x 16 (configurações 1A--1H
%   e 3A--3H) e outro para os circuitos 32 x 32 (configurações 2A--2H)."
% MOTIVO: os mapas por grupo (mapa_grupos_v4.py) escolhem o modo só com as configurações do
%   grupo (modo_fixo(p, occ_k, ks, ...) com ks = configurações do grupo). Os grupos 1 e 3
%   têm o mesmo circuito 16 x 16, mas podem receber modos diferentes. RESOLVIDO (2026-11-04):
%   a autora decidiu fazer a análise por grupo (modo e formatos escolhidos por grupo).
\begin{enumerate}
\item \textcolor{red}{o modo que deixa mais pares (configuração, condição de teste) com aumento de SNR de até 3~dB. Um par que não atinge $\mathrm{SER} = 10^{-2}$ conta como fora do limite;}
\item \textcolor{red}{a menor ocupação, pelo recurso mais usado na síntese;}
\item \textcolor{red}{a quantização \texttt{AP\_RND\_ZERO};}
\item \textcolor{red}{o menor aumento de SNR no pior par;}
\item \textcolor{red}{a ordem dos modos na lista.}
\end{enumerate}
% FONTE: mapa_cenarios_v4.py, modo_fixo (linhas 149-176): a chave é (pares dentro do
%   limite, ocupação, AP_RND_ZERO, pior perda, ordem dos modos). RESOLVIDO: AP_RND_ZERO
%   vem antes da pior perda (item acrescentado em vermelho).

\subsection{\texorpdfstring{\textcolor{red}{Mapas de aumento de SNR}}{Mapas de aumento de SNR}}\label{subsec:mapas-perda}
\textcolor{red}{O resultado da validação é mostrado em mapas de aumento de SNR, um para cada grupo de configurações. Cada linha do mapa é um formato, com o modo fixo do grupo. Cada coluna é uma configuração. A célula mostra o maior aumento de SNR entre as condições de teste. Os símbolos mostram as condições com aumento de SNR acima de 0,1~dB (vazados) e acima de 3~dB (cheios). Um x marca a configuração em que alguma condição não atinge $\mathrm{SER} = 10^{-2}$. Os mapas estão no capítulo de resultados.}
% FIGURAS: resultados/4_varredura_formatos/v4_formatos_juntos_16_configs/figuras/por_grupo/
%   mapa_perda_modo_fixo_ate_3dB_iter1_grupo{1,2,3}_modo_por_grupo.pdf (mapa_grupos_v4.py).
%   Grupos 2 e 3 já com as QPSK da v5. Escala de cor logarítmica, de 0,01 a 3 dB.

\subsection{\texorpdfstring{\textcolor{red}{Formatos escolhidos}}{Formatos escolhidos}}
\textcolor{red}{Os formatos foram escolhidos em cada grupo de configurações, com o modo fixo do grupo. Em cada grupo, foram escolhidos três formatos: o primeiro em que todas as configurações e condições atingem $\mathrm{SER} = 10^{-2}$; o primeiro em que todas atingem esse ponto com aumento de SNR de até 3~dB; e o primeiro em que todas atingem com aumento de até 0,1~dB. Os formatos são percorridos do menor para o maior $W$. A Tabela~\ref{tab:formatos-escolhidos} mostra os formatos escolhidos. Esses formatos foram implementados no Vivado, para medir a potência e as arquiteturas de cada condição de teste.}

\begin{table}[htbp]
\caption{\textcolor{red}{Formatos escolhidos em cada grupo de configurações.}}
\label{tab:formatos-escolhidos}
\centering
\footnotesize
\setlength{\tabcolsep}{3pt}
\begin{tabular}{>{\raggedright\arraybackslash}p{2.1cm}ccc}
\toprule
Critério & Grupo 1 & Grupo 2 & Grupo 3 \\
\midrule
Atinge $\mathrm{SER} = 10^{-2}$ & \begin{tabular}[c]{@{}c@{}}$\langle 22,12\rangle$\\ \texttt{AP\_RND\_ZERO}/\texttt{AP\_WRAP}\end{tabular} & \begin{tabular}[c]{@{}c@{}}$\langle 24,12\rangle$\\ \texttt{AP\_RND\_ZERO}/\texttt{AP\_WRAP}\end{tabular} & \begin{tabular}[c]{@{}c@{}}$\langle 24,12\rangle$\\ \texttt{AP\_RND\_ZERO}/\texttt{AP\_WRAP}\end{tabular} \\
\addlinespace
Aumento de até 3~dB & \begin{tabular}[c]{@{}c@{}}$\langle 24,12\rangle$\\ \texttt{AP\_RND\_ZERO}/\texttt{AP\_WRAP}\end{tabular} & \begin{tabular}[c]{@{}c@{}}$\langle 28,14\rangle$\\ \texttt{AP\_RND\_ZERO}/\texttt{AP\_WRAP}\end{tabular} & \begin{tabular}[c]{@{}c@{}}$\langle 24,12\rangle$\\ \texttt{AP\_RND\_ZERO}/\texttt{AP\_WRAP}\end{tabular} \\
\addlinespace
Aumento de até 0,1~dB & \begin{tabular}[c]{@{}c@{}}$\langle 28,14\rangle$\\ \texttt{AP\_RND\_ZERO}/\texttt{AP\_WRAP}\end{tabular} & \begin{tabular}[c]{@{}c@{}}$\langle 30,14\rangle$\\ \texttt{AP\_TRN}/\texttt{AP\_SAT}\end{tabular} & \begin{tabular}[c]{@{}c@{}}$\langle 30,14\rangle$\\ \texttt{AP\_TRN}/\texttt{AP\_SAT}\end{tabular} \\
\bottomrule
\end{tabular}
\legend{\textcolor{red}{Fonte: a autora. Abaixo do formato, o modo de quantização e o de estouro.}}
\end{table}
% FONTE: mesmo cálculo de mapa_grupos_v4.py (modo_fixo com as configurações do grupo, limite
%   3 dB, 1 iteração, v4 + v5). Script de conferência no scratchpad da sessão 981d36d9
%   (formatos_por_grupo.py). Pior perda no grupo (todas as configurações e condições):
%   G1: <22,12> 3,33; <24,12> 0,45; <26,12> 0,16; <28,12> 0,17; <28,14> 0,02; <30,14> 0,01.
%   G2: <22,12> x; <24,12> 4,12; <26,12> 3,76; <28,12> 3,67; <28,14> 0,34; <30,14> 0,04.
%   G3: <22,12> x; <24,12> 2,55; <26,12> 2,45; <28,12> 2,46; <28,14> 0,13; <30,14> 0,03.
%   Formatos de 6 e 8 bits inteiros têm x em todos os grupos.
% "Grupo 1/2/3": os nomes G1-G3 eram só internos (memória tcc-nomes-configuracoes). Aqui o
%   grupo é o número da configuração; CONFERIR se a autora quer "Grupo 1" na tese ou outro
%   rótulo (ex.: "1A--1H").
% <26,12>: com a análise por grupo, não é escolhido em nenhum grupo.
% ANTIGO (lista por todas as configurações juntas, antes da análise por grupo):
% Da validação, foram escolhidos os formatos para a implementação no Vivado, a medição de potência e as arquiteturas de cada condição de teste:
% \begin{itemize}
% \item $\langle 22,12\rangle$: o primeiro formato em que todas as configurações e condições atingem $\mathrm{SER} = 10^{-2}$;
% \item $\langle 24,12\rangle$: o primeiro formato em que todas atingem $\mathrm{SER} = 10^{-2}$ com perda de até 3~dB;
% \item $\langle 28,14\rangle$: o primeiro formato em que todas atingem $\mathrm{SER} = 10^{-2}$ com perda de até 0,1~dB;
% \item $\langle 26,12\rangle$: % PENDENTE: critério da escolha (decidido em 2026-10-31).
% \item $\langle 30,14\rangle$ em \texttt{AP\_TRN}/\texttt{AP\_SAT}: % PENDENTE: critério.
%       % ATENÇÃO (1 iteração): a única justificativa registrada era atender T = 4
%       % (0,04 dB na [5]). Com T = 1, <28,14> já fica até 0,1 dB. Na primeira versão,
%       % decidir se <30,14> entra e por quê; na versão com 4 iterações (se houver
%       % tempo), essa é a justificativa natural.
% \end{itemize}
% Os quatro primeiros usam \texttt{AP\_RND\_ZERO}/\texttt{AP\_WRAP}.
% % FONTE: software/testes/LEIAME.md, rodada [5] modo_fixo_3dB; mapa iter1 (decisão de
% %   2026-11-04: o TCC usa só 1 iteração). "Primeiro" = menor W na ordem dos 10 formatos.
%
